\documentclass[11pt]{article}
\usepackage[margin=1in]{geometry}
\usepackage[T1]{fontenc}
\usepackage{lmodern}
\usepackage{microtype}
\usepackage{amsmath,amssymb}
\usepackage[colorlinks=true,linkcolor=blue,urlcolor=blue,citecolor=blue]{hyperref}
\setlength{\parindent}{0pt}
\setlength{\parskip}{0.7em}

\title{Nontrivial Case-3 Spin--Boson Renormalization\\
\large An explicit literature answer to arXiv:2508.00803}
\author{Literature-status note by \texttt{agent\_lane\_06}}
\date{22 July 2026}

\begin{document}
\maketitle

\textbf{Status.} \texttt{literature\_already\_answered}.  This is a
literature identification, not a new mathematical result of the run.

\section*{Source question}

Alvarez, Lill, Lonigro, and Valent\'in Mart\'in study generalized spin--boson
Hamiltonians and divide ultraviolet singularities into three cases.  Case~3
is the supercritical regime, in which at least one form factor $f$ satisfies
$f/\omega\notin L^2$.  In Section~1.2, PDF page~3, after describing their
triviality result, they write:

\begin{quote}
``The question whether it is possible to obtain a nontrivial Case 3
renormalization of certain GSB models remains open.''
\end{quote}

This is an existential question: one nontrivially renormalizable class in the
supercritical regime gives an affirmative answer.

\section*{Supporting answer}

Falconi, Hinrichs, and Valent\'in Mart\'in develop a wave-function
renormalization on a renormalized, generally non-Fock Hilbert space.  Their
Theorem~3.18, ``Renormalization of the SB model'' (PDF page~22), applies to a
bounded normal spin interaction operator $B$ and every distributional source
$v\in\mathcal T'$.  It constructs a renormalized Hamiltonian
$H_{\mathrm{ren}}(B,v)$ and proves:

\begin{itemize}
  \item self-adjointness on the dressed domain, with a specified core;
  \item a uniform lower bound by $-\lVert A\rVert$; and
  \item generalized strong resolvent convergence of regularized Hamiltonians
        as $v_\alpha\to v$ (and generalized norm resolvent convergence in the
        stated Hilbert-space subcase).
\end{itemize}

Because $v$ is arbitrary in $\mathcal T'$, the theorem includes sources with
$v/\omega\notin\mathfrak h$, precisely the supercritical regime excluded from
energy-only renormalization on the original Fock space.  The construction is
therefore a nontrivial Case-3 renormalization and answers the source question
affirmatively.

\section*{Provenance and scope}

The identification is explicit in the literature.  The supporting paper
cites arXiv:2508.00803 in its account of critical and supercritical
spin--boson renormalization.  It supersedes the withdrawn arXiv:2508.00805,
whose title and abstract explicitly announced that its construction solved
the triviality problem for unitarily renormalized supercritical spin--boson
models.

The result requires $B$ to be bounded and normal.  Thus, it does not provide a
general theory for arbitrary nonnormal interaction matrices or arbitrary
multi-interaction models.  Those restrictions do not affect the affirmative
answer to the source's existential question.

\section*{Search evidence}

The source bibliography already points to arXiv:2508.00805.  Inspection of its
current arXiv record showed that it was withdrawn in March 2026 and superseded
by arXiv:2603.07045.  The replacement paper's abstract, Section~3 overview,
Theorem~3.18, and citation to arXiv:2508.00803 establish the answer and its
scope.

\begin{thebibliography}{9}
\bibitem{source}
B. Alvarez, S. Lill, D. Lonigro, and J. Valent\'in Mart\'in,
\emph{Renormalization of generalized spin--boson models with critical
ultraviolet divergences}, arXiv:2508.00803v2 (2025),
\url{https://arxiv.org/abs/2508.00803}.

\bibitem{answer}
M. Falconi, B. Hinrichs, and J. Valent\'in Mart\'in,
\emph{Wave Function Renormalization for Particle-Field Interactions},
arXiv:2603.07045v1 (2026),
\url{https://arxiv.org/abs/2603.07045}.

\bibitem{superseded}
M. Falconi, B. Hinrichs, and J. Valent\'in Mart\'in,
\emph{Non-Trivial Renormalization of Spin-Boson Models with Supercritical
Form Factors}, withdrawn arXiv:2508.00805 (2025--2026), superseded by
\cite{answer}.
\end{thebibliography}

\end{document}
